\begin{center}
    {\LARGE \textbf{2008无机化学}} \\[2ex]
    {\large 时量180分钟 \quad 满分100分}
\end{center}

\vspace{0.5cm}
\noindent\textbf{注意事项:}
\begin{enumerate}[label=\arabic*., parsep=0pt, itemsep=0pt]
    \item 答题前,考生先将自己的姓名、学号、考场号、座位号填写在答题卡上,将条形码准确粘贴在考生信息条形码粘贴区。
    \item 考试结束后,将本试卷与答题纸一并收回。
    \item 回答各个问题时,将答案写在答题纸上,写在本试卷上无效。
\end{enumerate}

\vspace{0.5cm}
\noindent\textbf{一、完成并配平下列化学反应方程式（24 pts.）}

\begin{enumerate}[label=\arabic*., start=1, itemsep=1.5ex]
    \item \ce{Pt + HNO3 + HCl =}
    \item \ce{BiO3^{-} + Mn^{2+} + H^{+} =}
    \item \ce{SO2Cl2 + H2O =}
    \item 将\ce{H2}通入灼热的\ce{WO3}
    \item 氢碘酸在常温下被空气中的氧气所氧化
    \item \ce{Al2O3}、\ce{Cl2}和碳共热
\end{enumerate}

\vspace{0.5cm}
\noindent\textbf{二、合成与制备（22 pts.）}

\begin{enumerate}[label=\arabic*., start=1, itemsep=1.5ex]
    \item 以单质硫和碳酸钠（\ce{Na2CO3}）为原料制备硫代硫酸钠。
    \item 试说明工业上以\ce{Ca3(PO4)2}为主要原料制取单质磷的方法。
    \item 试说明以软锰矿（\ce{MnO2}）为原料生产\ce{KMnO4}的过程。
\end{enumerate}

\vspace{0.5cm}
\noindent\textbf{三、鉴别与分离（28 pts.）}

\begin{enumerate}[label=\arabic*., start=1, itemsep=1.5ex]
    \item 黄色固体样品，可能是\ce{PbCrO4}或\ce{BaCrO4}，试加以鉴别。
    \item 白色固体样品，可能是氯化汞或氯化亚汞，试加以鉴别。
    \item 一无色溶液，可能是\ce{Na2CO3}或\ce{Na2SiO3}，试加以鉴别。
    \item 试设计方案对混合溶液中的离子进行分离：\ce{Ag^{+}}、\ce{Al^{3+}}、\ce{Cr^{3+}}、\ce{Sb^{3+}}、\ce{K^{+}}
\end{enumerate}

\vspace{0.5cm}
\noindent\textbf{四、简要回答下列问题（42 pts.）}

\begin{enumerate}[label=\arabic*., start=1, itemsep=1.5ex]
    \item 试画图表示链四聚和环四聚磷酸分子内原子间的键联关系，并写出链$n$聚和环$n$聚磷酸分子的通式。
    \item 试写出一氯·三氨·二水合钴(III)离子的化学式，它有多少种空间异构体？试全部画出。（以$\ce{Ma1b2c3}$形式表示）
    \item 试用价层电子对互斥理论判断下列分子或离子的空间构型，并说明其中心原子的杂化方式：
    
    {\centering
    \ce{PCl5}，\ce{SF6}，\ce{H3O^{+}}，\ce{NO2}，\ce{SF4}，\ce{XeOF4} \par
    }
    \item 闪锌矿（\ce{ZnS}）、辰砂（\ce{HgS}）在空气中焙烧的主要产物各是什么物质？试分析其不同的原因。
    \item 某原电池的一个电极由锌片插到$0.10\ \mathrm{mol\cdot dm^{-3}}$\ \ce{ZnSO4}溶液中构成；另一个电极由锌片插到混合溶液中构成，该混合溶液中$[\ce{Zn(NH3)4^{2+}}]=0.10\ \mathrm{mol\cdot dm^{-3}}$；$[\ce{NH3}]=1.0\ \mathrm{mol\cdot dm^{-3}}$。测得原电池的电动势为$0.278\ \mathrm{V}$，试求\ce{Zn(NH3)4^{2+}}的$K_{\text{稳}}$。
    \item 副族元素的第一电离能$I_1/\mathrm{kJ\cdot mol^{-1}}$数据如下：
    \[
    \begin{array}{ccccccccc}
    \ce{Ti} & \ce{V} & \ce{Cr} & \ce{Mn} & \ce{Fe} & \ce{Co} & \ce{Ni} & \ce{Cu} & \ce{Zn} \\
    659 & 651 & 653 & 717 & 762 & 760 & 737 & 746 & 906 \\
    \ce{Zr} & \ce{Nb} & \ce{Mo} & \ce{Tc} & \ce{Ru} & \ce{Rh} & \ce{Pd} & \ce{Ag} & \ce{Cd} \\
    640 & 652 & 684 & 702 & 710 & 720 & 804 & 731 & 868 \\
    \ce{Hf} & \ce{Ta} & \ce{W} & \ce{Re} & \ce{Os} & \ce{Ir} & \ce{Pt} & \ce{Au} & \ce{Hg} \\
    659 & 728 & 759 & 756 & 814 & 865 & 864 & 890 & 1007 \\
    \end{array}
    \]
    \begin{enumerate}[label=(\alph*), itemsep=1ex]
        \item 试指出影响同族元素第一电离能大小的因素；
        \item 比较三个过渡系列元素的第一电离能，试指出三者之间的关系，并说明原因。
    \end{enumerate}
\end{enumerate}

\vspace{0.5cm}
\noindent\textbf{五、推理判断（34 pts.）}

\begin{enumerate}[label=\arabic*., start=1, itemsep=1.5ex]
    \item 黑色化合物（A）在高温下分解成（B）和（C）。红色物质（B）与稀硫酸反应沉淀出（D），浓缩溶液时析出（E）。红色物质（D）不与稀盐酸反应，但可以与热浓盐酸作用放出气体（F）。蓝色物质（E）加热时可以变成（G），白色物质（G）与氢氧化钠溶液反应生成蓝色沉淀（H），（H）溶于氨水得深蓝色溶液，说明有（I）生成。物质（D）在无色气体（C）中灼烧生成化合物（A）。使过量气体（F）通过灼热的（A），反应充分时生成（D），若气体不充足时产物中有（B）。
    试给出（A）、（B）、（C）、（D）、（E）、（F）、（G）、（H）和（I）的化学式。
    
    \item 红色化合物（A）经稀硝酸处理后得黑色沉淀物（B）和无色透明溶液，从该溶液中可以结晶出无色晶体（C）。晶体（C）加热时放出红棕色气体（D）。（B）在酸性介质中与二价锰离子作用，溶液逐渐变为紫红色，说明有化合物（E）生成。若向（C）的水溶液中滴稀盐酸有白色沉淀（F）生成，（F）不溶于氨水；若向（C）的水溶液中滴入碘化钾溶液，则生成黄色沉淀（G）；若向（C）的水溶液中滴入氢氧化钠溶液，有白色沉淀（H）生成，氢氧化钠过量则沉淀消失。
    试给出（A）、（B）、（C）、（D）、（E）、（F）、（G）和（H）的化学式。
\end{enumerate}